\documentclass[11pt,a4paper]{article}

% --- Packages ---
\usepackage[utf8]{inputenc}
\usepackage[english]{babel}
\usepackage{booktabs}   % For professional tables
\usepackage{geometry}   % For page margins
\usepackage{hyperref}  % For clickable links and citations
\usepackage{natbib}    % For bibliography management
\usepackage{listings}
\usepackage{color}
\usepackage{graphicx} %package to manage images
\graphicspath{ {./Assets/} }

\definecolor{dkgreen}{rgb}{0,0.6,0}
\definecolor{gray}{rgb}{0.5,0.5,0.5}
\definecolor{mauve}{rgb}{0.58,0,0.82}

\lstset{frame=tb,
  aboveskip=3mm,
  belowskip=3mm,
  showstringspaces=false,
  columns=flexible,
  basicstyle={\small\ttfamily},
  numbers=none,
  numberstyle=\tiny\color{gray},
  keywordstyle=\color{blue},
  commentstyle=\color{dkgreen},
  stringstyle=\color{mauve},
  breaklines=true,
  breakatwhitespace=true,
  tabsize=3
}

\geometry{margin=1in}
\hypersetup{
    colorlinks=true,
    linkcolor=blue,
    citecolor=blue,
    urlcolor=blue
}

% --- Document Metadata ---
\title{A Comparative Analysis of Cost-Effective Volunteer Recruitment Strategies in the UK}
\author{Brandon Antony Wright}
\date{\today}

\begin{document}

\maketitle

\begin{abstract}
A small report on the technological methodologies to support volunteering charities in the modern world. 
\end{abstract}

\section{Introduction}

Following a recent directive outlines the requirements for a proposed regional database, this report addresses the strategic foundation needed to launch a self-sustaining voluntary sector platform. The initial brief specified a system capable of delivering a voluntary sector "job" or "skills" match, incorporating career development tracking, user registration, a reporting facility, and a monetization model to fund ongoing management, all while strictly avoiding GDPR compliance issues. Please note that a full in depth look into the research done can be found \href{https://github.com/Barnold8/voluntary-sector-jobmatch/blob/main/Docs/Email.md}{here}, also note it's more a stream of consciousness rather than a professional piece of work

\hfill \break
Specifically, this report will look at three cost effective measures based on research looking into localised and national charities and how they operate. 

\section{Methodology}

To begin this process, I looked at more generalised charities that are popular in the UK. I looked into how they operate and speculated on their cost effective measures. This was done to start the groundwork of my understanding on the charity space and I have acquired a lot more understanding over time. However, it should be expected that I am wrong in some points since this is a week project and it takes longer to summarize a depth of knowledge on any sector. 

\begin{itemize}
    \item \href{https://reachvolunteering.org.uk/}{Reach} 
    \item \href{https://www.charityjob.co.uk/?source=nav}{CharityJob}
    \item \href{https://www.doit.life}{Do It}
\end{itemize}

My summary on these charities is that they are like a megaphone to smaller charities around the nation. Bolstering the volunteering sector into posting opportunities and pushing volunteers forward to bigger and better things through volunteer work. I would personally recommend looking into these charities and contacting them to collaborate where appropriate. 

\hfill \break
After the national charities, I looked at more localised charities around the UK looking into wealthier areas and poorer areas according to the 2021 sensus data (Household deprivation - census maps, ons) to gain an understanding of the different ways to obtain volunteers through differing price points, understanding how a charity thrives in a low cost environment comparative to a high cost environment.

\hfill \break
\hfill \break

These are the charities I looked into for my research: 

\begin{itemize}
    \item \href{https://togetherco.org.uk/}{ TogetherCo - Brighton}
    \item \href{https://manchestercommunitycentral.org/volunteer-centre-manchester}{Manchester Community Central - Manchester} 
    \item \href{https://wtonvolunteers.org/}{Wtonvolunteers - Wolverhampton} 
    \item \href{https://www.thekidsnetwork.org.uk/i-want-to-mentor?gad_source=1&gad_campaignid=22790719866&gbraid=0AAAAADRkgo9ZgrjOC3w8BTZU-XkR_5Yrm&gclid=Cj0KCQjw_vnQBhCxARIsADcZyxJlH0f1fgNgS3j_ExdetIzvqbD3oBxhusiLByw8EieP9HI863fcQbQaAlPtEALw_wcB}{The Kids Network - London} 
    \item \href{https://volunteercentrenewcastle.org.uk/}{Volunteer Centre Newcastle - Newcastle} 
    \item \href{https://vast.org.uk/}{Vast - Stoke on Trent} 
    \item \href{https://www.bvsc.org/}{BVSC - Birmingham} 
    \item \href{https://www.rainbows.co.uk/}{Rainbows - East Midlands} 
\end{itemize}

\section{Analysis / Findings}

All the charities I researched have some form of registration for their volunteers and the organisations who wish to post volunteer roles. This will be important later to consider what tech will be required to facilitate registration methodologies if any at all. On top of this, all charities had some role searching capabilities to help volunteers find the roles they are looking for. A majority of these charities use some s

 of form system to gather user/organisation data to register them with rather than the typical registration page. When considering these forms and their usage, we need to consider how the form is being handled when it is submitted. Is the form data appended to a row in a CSV? Is it logged in a database? Is it just sent via E-Mail to an employee? A cost effective strategy found specifically in the Newcastle charity researched, was the usage of Office 365. Using the forms software, employees can generate useful forms and be assured that its GDPR compliant in terms of securely passing information over the internet using military grade encryption and allowing users to delete data at request. 

\hfill \break
What was also found was the newsletter and/or recent news sections. While not the top priority for a charities success, it can be a good idea to show what is happening in the surrounding area for volunteer work. In the case of the newsletter, its a great autonomous reminder to registered volunteers about the charity, keeping said charity in the back of their mind for longer. It should also be noted that 3rd party platforms already exist to handle newsletters automatically for organisations. The following were found 

\begin{itemize}
    \item \href{https://www.cleverreach.com/en/newsletter-tool/creating-and-designing-newsletters/?tl=o8ejlo-aw_uk_cn&gad_source=1&gad_campaignid=1370877092&gbraid=0AAAAADj-Y5p11HESNGlfDIL4UX7zsALFv&gclid=CjwKCAjwxITRBhBYEiwA6mZm7ZlYPj1hoQf-Fzp6Ay-qqrejPVhJTuu8OYbZEhog0A4dmfQC6n6xlBoCZxUQAvD_BwE}{CleverReach}
    \item \href{https://exclaimer.com/lp/14-day-free-trial/?hstk_creative=810959687171&hstk_campaign=23892133625&hstk_network=googleAds&creative=810959687171&keyword=email%20brand%20consistency&matchtype=b&network=g&device=c&utm_term=email%20brand%20consistency&utm_campaign=UK%20%7C%20Search%20%7C%20Generic%20%7C%20Broad%20Match%20A%2FB%2FC%20Landing%20Page%20Test&utm_source=adwords-search&utm_medium=cpc&utm_content=Broad%20-%20Brand%20Identity&hsa_acc=4500841410&hsa_cam=23892133625&hsa_grp=195322110445&hsa_ad=810959687171&hsa_src=g&hsa_tgt=kwd-2424627345829&hsa_kw=email%20brand%20consistency&hsa_mt=b&hsa_net=adwords&hsa_ver=3&gad_source=1&gad_campaignid=23892133625&gbraid=0AAAAADzh7UANuNVAaIuF3fxvl9toXU9L-&gclid=CjwKCAjwxITRBhBYEiwA6mZm7UwT0Ct1B1Qu8LeeRH96L80FVK0YAwWXM3-chf2Cw3-fzRebgScf5hoCSUwQAvD_BwE}{Exclaimer}
    \item \href{https://monday.com/crm/lp/campaigns?cq_src=google_ads&cq_cmp=23044226584&cq_term=email%20marketing%20software&cq_plac=&cq_net=g&cq_plt=gp&utm_medium=cpc&utm_source=adwordssearch&utm_campaign=gb-en-prm-workos-monday_campaigns-email_marketing-h-search-desktop-core-aw&utm_keyword=email%20marketing%20software&utm_match_type=p&cluster=crm_industries&subcluster=crm_healthcare&ati=&utm_adgroup=email%20marketing&utm_banner=775017905004&gad_source=1&gad_campaignid=23044226584&gbraid=0AAAAADeiQJuzMvjpNP5Fb8wJ0IXbxiNiK&gclid=CjwKCAjwxITRBhBYEiwA6mZm7Z7wvFhfbakc2QlzwoKGuzz8tRepHgZ_hfrbt0brv2AbROvnt8Z1aRoCZQ0QAvD_BwE}{Monday Campaigns}

\end{itemize}

It should also be noted that an alternative to these paid subscription services is quite simple and cost effective. Simply writing an E-Mail and using the BCC feature to all registered newsletter participants works just as well if not taking more time than a paid subscription.  

\hfill \break
Booking a meeting with a charity was found to also be a common methodology. What was found was in all cases where a booking was possible, \href{https://calendly.com/}{calendly} was used. The likely reason booking is common is because its a great alternative to rigid account registration where empathy and complex triaging is none existent.  


\section{The three cost measures}

This section is here to suggest three methodologies. The low cost will assume a charity with next to no funding. The medium cost will most closely resemble the average charity, not struggling but not prospering either. The high cost will assume a charity with a lot of funding to put wherever it is needed. The assumption for this section is that some website already exists for the charity that is accessible to be edited as needed and is already running on a web server somewhere on the internet, this fact will not be considered as one of the cost measures. The requested measures from the client's E-Mail will be found in the generalised "Storing user and organisation accounts while also displaying job postings" chapter in each cost section. The reasoning behind this because all the main requirements are closely coupled with that problem as it stands. 

\subsection{Low Cost}

The low cost approach for this needs to pay as little money as possible while providing the best system and experience to organisations and volunteers.

\subsubsection{Storing user and organisation accounts while also displaying job postings}

Now storing and displaying volunteering roles can work one of two ways here. The first is completely free but time consuming whereas the second is cheap but automatic. 

\hfill \break
The first would be to have a form for volunteers and organisations. The volunteer form would take specific information like their name, address, skill set, experience and if they are aiming to further their career or not. The organisation form would be similar but require respective fields like the address of the role, days that the role takes place, the time frame of the role, the hours for the days of the role etc. In this, an employee will have to contact and have meetings with volunteers to understand what they are looking for and follow up on possible volunteer roles that suit them. The employee must be trained in GDPR guidelines to ensure that no sensitive data is kept longer than is necessary. So while its free, it may be more of a headache and cost more in employee fees but something that \textbf{can} be considered. In terms of reporting statistics for this section, it may just be as simple as scrolling down the spreadsheet file and seeing what row number has data in it. For an updating value that can stay at the top you can use the following: 
\hfill \break
\begin{lstlisting}
=COUNTA(A:A)

Note: pick a column letter with data in it, this will update as the sheet updates
\end{lstlisting}

\hfill \break
The second would be using a Volunteer Management System like \href{https://www.paamapp.com/}{PAAM} which can be embedded in a webpage on the website itself in an \href{https://en.wikipedia.org/wiki/HTML_element#Frames}{i-frame}. This can act as a pseudo database allowing role postings to be somewhat automated and also allows volunteers to register too. It can also skill match volunteer profiles to volunteer roles too (Bonterra, 2025). Due to this being a third party, the charity themselves do not have to worry about all the complex details that are required to keep the system running. The most the charity will have to do is embed the system within a webpage. With the skill matching system in place in the VMS, it also bolsters career development (Evans, 2024)! The only assumed downside to a VMS is the admin behind it all. The tech side of the system is all handled but the admin of it most likely isn't. It should also be noted that the VMS you choose may or may not have capabilities to enforce paid registration for organisational accounts. It may require some money to pay a developer to edit the current existing site to make this paid registration, prices may vary. 

\subsection{Medium Cost}

The medium cost doesn't concern itself with cutting costs at every possible turn but it also doesn't want to spend as much money as it has. 

\hfill \break

In this section we are going to use the \href{https://n8n.io/}{n8n} platform. N8n is an automation platform that can facilitate most of the requirements of the client while being customisable to meet new needs and is a medium cost effective measure. The downside to this approach is that a knowledgable person will need to set it up and will be required if changes are to be made to the system. With this we can create our own forms like the following:
\begin{center}
    \includegraphics[width=300]{Assets/n8n_form.png}
\end{center}
And with this, we can have the form be a link on the already existing website, maybe its a registration button or a hyperlink. On the back end, when the user submits the form, it will go through a system like the following:
\begin{center}
    \includegraphics[width=500]{Assets/n8n_automation.png}
\end{center}

This \textbf{very} simple example shows a system that appends the new registered user to a CSV and then send that email to an admin. The email step has an x because at testing it required some technical step that I couldn't set up at the time. All of the nodes in that graph are customisable too, which means we can sift through data, categorise it, report it back to an admin via email etc. The only downside to n8n that I am aware of in respect to the client's requirements is a proper data visualisation methodology like seeing user metrics on a user friendly dashboard. The n8n workflow can automatically add and remove from a database allowing for automatic sensitive data removal and appending new roles to a database that the website uses to list it's volunteer roles. Unfortunately n8n does not have any system in place to view and interact with data since its primarily a automation platform. But this database can be directly linked with the usage of the website using whatever frameworks are already in place with the website or implementing a new one. One of the great things with n8n in particular is the integrated payment platforms that come out of the box, you just add a payment related node and once a payment is authenticated, you can pass the process along to the next step. This is what we can do for the organisation side of things, requiring an organisation to pay a one time fee to register with the service or for every post (however the second option may put a lot of organisations off). Because the system is automated, we can add a procedure to check over the data for us and purge any data that would violate GDPR guidelines just before they would violate them. Again, with an automated platform we can set intervals for skill matching with organisations to occur. The system would every so few days or weeks compare the entire set of users and their skills with the current volunteer roles available and match the as best as possible. And with the AI integration, we can use the intelligence side of AI to match candidates better to roles than a predertimed algorithm could. 

\hfill \break
So in short, this option becomes like the low cost version, but with more reporting, more refined data analysis techniques and no headaches due to automated system features doing all the heavy lifting. 

\subsection{High Cost}

The high cost approach for this needs doesn't concern itself with how much things cost because we can assume a charity somehow has funds to spare on development costs like this. So how do we use money to make the most robust and diverse system for a volunteering charity? Since we are having this made from scratch, we are going to have to consider things that wouldn't need to be considered when using 3rd party systems. For example, the resilience of the system itself, can it handle large droves of users? Can it handle cyber attacks from criminals? This section will have a lot of technical jargon but I will try my best to explain it in laimens terms such that anybody can understand irregardless of experience level. 

The system I have devised can be seen in the diagram below:

\begin{center}
    \includegraphics[width=500]{Assets/High-cost-strategy.png}
\end{center}

The "client app / Web front end" is the website itself. In this part we will handle basic security dilemmas that can happen from attackers. Please note the website itself is just what we send to a client when they go to the URL for the website. The website itself would include a search page for volunteers to search for a role, registration page, so organisations and volunteers can register with the charity, an about us section among many of the other pages you see on the average charity web page. Because this is custom, we can use whatever frameworks we like so we can make the website look to the clients specification. With the front end being whatever we want it to be, we can use whatever payment platforms we want and change them out if needed. We could use Paypal, Stripe among other payment platforms. This means that we can use a subscription based model or a one time payment model for organisations to pay to be able to post to the charity's role post section. This means that we will have no spam and the charity can keep up with management costs. With this too, we can also implement any advertisement platform we want too to increase funds to the charity through ad sense. 

\hfill \break

The Load balancer is a common idea in server programming to distribute the web traffic coming into the servers evenly to avoid overloading one server in a network. You can think of this like not having too many people going through one ticket gate at a train station but being evenly distributed through all of them such that the gates don't fail, its the same premise. 

\hfill \break

The web servers are what the browser actually communicates with to provide the services the volunteers and organisations require. On top of this, they'll come with automated software that'll automatically purge any data that voids GDPR guidelines and sends user count reports among other statistics to the admin when they request it on the web front end. 

\hfill \break

The relational database is where all the information we need to keep will be stored. This will be stuff like, organisational accounts, user accounts, role postings, newsletter updates etc. Because the database is relational it means pieces of data are linked with others. This means that we can match volunteers skills that they have included in their profile with roles that have been posted and we can do this at any interval we please. The actual processing of this could either be done on a web server itself, or to reduce any stress on the webserver, this could be hosted on a separate machine or on the cloud as a microservice. 

\hfill \break
The reason this costs the most isn't necessarily the initial development costs but the upkeep and running of multiple servers and possible maintenance that would occur over time.  


\section{Conclusion}

The conclusion to the charity technology problem is mostly solved already with the VMS. However, there are aches and pains still relevant to this day with bad website designs and a VMS doesn't allow for customisability for systems. However it does provide all of the solutions to the clients problems so it's what I would suggest for simplicity and cost sakes. Personally, I would use the medium cost effective measure using the automation platform mentioned. With this you can still keep the old website but link to the nodes from the automation platform allowing for automatic systems to fire up when a user registers, when a new payment comes in from an organisation. The medium price point would provide all the proper statistical data to analyse how to properly move the charity forward and provide insights that might not be known with the possibly limited VMS. However, if the charity ever becomes way too big that an automation platform can not suffice, the high cost solution may be the best outcome due to the unlimited potential for the client's requirements. The high cost will be overkill for even the biggest charities out there irregardless since many 3rd party platforms are doing all of the heavy lifting anyway.  

\section{Less technical but possibly valuable outlook}

When it comes to the modern world, it matters to keep up with rising trends and how the youth and the general public view charities and the voluntary sector in general. With social media being as popular as it ever has been, and now how so many people are stuck on it day to day (Number of social media users worldwide 2030 2026), it may be worth using these platforms to gain attention and steer more eyes to the voluntary sector (R/digitalmarketing on reddit: How can businesses use social media marketing). This could increase funding, generalised interest and possibly bring people to the charity who have never heard of it before. Personally, I had never heard of *REDACTED* until it was brought up through word of mouth. Even in todays age, it still somewhat comes down not to what you know but who you know. It may even be more cost effective to hire a social media manager to post on social media, follow trends and garner positive attention for the charity (Why a strong social media presence is essential for business growth in 2026). Another methodology to gain new volunteers would be paying for ads. A simple example would be paying for the sponsored result on google search results, this way you can appear at the \textbf{very} top of search results. A way to earn funds, that may be in distaste to a charity, im not sure, would be to physically have ads on the charity website. That way you can turn traction to the website to direct funding to the charity. 

% --- References Section ---
\begin{thebibliography}{9}

\bibitem{Blackcountrychamber}

Why a strong social media presence is essential for business growth in 2026 (no date) Why a Strong Social Media Presence Is Essential for Business Growth in. Available at: https://www.blackcountrychamber.co.uk/news/article/why-a-strong-social-media-presence-is-essential-for-business-growth-in-2026/ (Accessed: 09 June 2026). 

\textit{blackcountrychamber}

\bibitem{Bonterra}

Bonterra (2025) 19 best volunteer management solutions for nonprofits, Bonterra. Available at: https://www.bonterratech.com/blog/volunteer-management-software (Accessed: 09 June 2026). 

\textit{Bonterra (2025) }.

\bibitem{census2021}
Household deprivation - census maps, ons (no date) Home - Office for National Statistics. Available at: https://www.ons.gov.uk/census/maps/choropleth/population/household-deprivation/hh-deprivation/household-is-deprived-in-one-dimension (Accessed: 09 June 2026). 
\textit{Census 2021 Data}.

\bibitem{Evans, A.}

Evans, A. (2024) Is volunteer management important?, TeamKinetic. Available at: https://teamkinetic.co.uk/blog/2024/03/is-volunteer-management-important/ (Accessed: 09 June 2026). 

\textit{Evans, A }

\bibitem{R/digitalmarketing}

(No date) R/digitalmarketing on reddit: How can businesses use social media marketing. Available at: \href{https://www.reddit.com/r/DigitalMarketing/comments/t8vhqt/how_can_businesses_use_social_media_marketing/}{R/digitalmarketing} (Accessed: 09 June 2026). 

\textit{R/digitalmarketing }

\bibitem{Statista}

Published by  Statista Research Department and 16, M. (2026) Number of social media users worldwide 2030, Statista. Available at: https://www.statista.com/forecasts/278414/number-of-worldwide-social-network-users/ (Accessed: 09 June 2026). 

\textit{Statista}




\end{thebibliography}

\end{document}